\chapter{\abstractname}

%TODO: Abstract
This thesis analysis common implementations of sampling algorithms used to generate variates of the Zipfian distribution, with special attention given
to the use cases found in storage systems and database technology research. Since the late eighties, there has been increased interest in evaluating the performance
of those systems under very skewed workloads, which are often modeled using the Zipfian distribution. While the theory behind some sampling techniques for the Zipfian distribution is well grounded, 
practical implementations vary in their accuracy and performance, which can lead to incorrect results in the evaluation of storage systems. This work presents a comprehensive analysis of the most common sampling algorithms 
found in benchmarking suites across academia and industry, and offers a comparison in terms of precision, speed and applicability to the aforementioned research areas. Parallely, two programs were written to help analysis: a benchmark tool was developed to run experiments on different implementations and
implementation of different variate generators was done in the programming language C++. Evaluation shows that most found implementations are suitable for the use cases in storage systems and database technology research, but some, such as FIO, are to be avoided. 
It was also revealed that table lookups might be an interesting alternative to established method, but further investigation is required.